% Bagian: turing-machines
% Bab: machines-computations
% Subbagian: combining-machines

\documentclass[../../../include/open-logic-section]{subfiles}

\begin{document}

\olfileid[id]{tur}{mac}{cmb}
\olsection{Menggabungkan Mesin Turing}

\begin{explain}
Contoh-contoh mesin Turing yang telah kita lihat sejauh ini cukup sederhana.
Namun, sesungguhnya setiap masalah yang dapat diselesaikan dengan bahasa
pemrograman modern apa pun juga dapat diselesaikan dengan mesin Turing. Untuk
membangun mesin Turing yang lebih rumit, penting bagi kita untuk memastikan
bahwa mesin-mesin itu dapat digabungkan. Dengan demikian, kita dapat membangun
mesin untuk menyelesaikan masalah yang lebih rumit dengan memecah prosedurnya
menjadi bagian-bagian yang lebih sederhana. Jika kita dapat menemukan cara
alami untuk menguraikan masalah rumit menjadi bagian-bagian penyusunnya, kita
dapat menangani masalah itu dalam beberapa tahap: membuat beberapa mesin
Turing sederhana lalu menggabungkannya menjadi satu mesin yang dapat
menyelesaikan masalah tersebut. Hal ini sangat penting ketika kita membahas
Masalah Penghentian pada bagian berikutnya.

Bagaimana kita menggabungkan mesin Turing $M = \tuple{Q, \Sigma, q_0, \delta}$
dan~$M' = \tuple{Q', \Sigma', q_0', \delta'}$? Kita gunakan konfigurasi pita
setelah $M$ berhenti sebagai konfigurasi masukan bagi jalannya mesin~$M'$.
Untuk memperoleh satu mesin Turing $M \frown M'$ yang melakukan hal tersebut,
lakukan langkah-langkah berikut:
\begin{enumerate}
    \item Nomori ulang (atau beri label ulang) semua keadaan~$Q'$ dari~$M'$
    agar $M$ dan~$M'$ tidak memiliki keadaan yang sama
    ($Q \cap Q' = \emptyset$).
    \item Keadaan-keadaan $M \frown M'$ adalah $Q \cup Q'$.
    \item Alfabet pitanya adalah $\Sigma \cup \Sigma'$.
    \item Keadaan awalnya adalah~$q_0$.
    \item Fungsi transisinya ialah fungsi $\delta''$ yang diberikan oleh:
    \[\delta''(q,\sigma) =
    \begin{cases}
      \delta(q,\sigma) & \text{jika $q \in Q$ dan $\delta(q,\sigma)$ terdefinisi}\\
      \delta'(q,\sigma) & \text{jika $q \in Q'$}\\
      \tuple{q_0', \sigma, \TMstay} & \text{jika $q \in Q$ dan
      $\delta(q,\sigma)$ tidak terdefinisi.}
    \end{cases}\]
\end{enumerate}
Mesin yang dihasilkan menggunakan instruksi-instruksi~$M$ ketika berada dalam
keadaan $q \in Q$, dan instruksi-instruksi~$M'$ ketika berada dalam keadaan
$q \in Q'$. Ketika mesin berada dalam keadaan $q \in Q$ dan memindai simbol
$\sigma$ yang tidak memiliki transisi dalam $M$ (yakni ketika $M$ semestinya
berhenti), mesin memasuki keadaan awal~$M'$ tanpa mengubah isi pita maupun
posisi kepala.

Perhatikan bahwa kecuali mesin~$M$ bersifat disiplin, kita tidak mengetahui
letak kepala pita ketika $M$ berhenti. Karena itu, pada konfigurasi akhir~$M$,
kepala belum tentu memindai petak~$1$. Hal ini penting diingat ketika
menggabungkan mesin.
\end{explain}

\begin{ex}
\emph{Menggabungkan Mesin:} Kita akan merancang mesin yang, ketika dijalankan
pada masukan berupa dua blok~$\TMstroke$ dengan panjang $n$ dan~$m$, berhenti
dengan satu blok yang berisi $2(m+n)$ buah~$\TMstroke$ pada pita. Untuk
membangun mesin ini, kita dapat menggabungkan dua mesin yang sudah dikenal:
mesin penjumlahan dan mesin pengganda. Kita mulai dengan menggambar diagram
keadaan mesin penjumlahan.
\[
\begin{tikzpicture}[->,>=stealth',shorten >=1pt,auto,node distance=2.8cm,
                    semithick]
  \tikzstyle{every state}=[fill=none,draw=black,text=black]

  \node[initial,state] (A)              {$q_0$};
  \node[state]         (B) [right of=A] {$q_1$};
  \node[state]         (C) [right of=B] {$q_2$};

  \path (A) edge node {\TMtrans{\TMblank}{\TMstroke}{\TMstay}} (B)
            edge [loop above] node {\TMtrans{\TMstroke}{\TMstroke}{\TMright}} (A)
        (B) edge [loop above] node {\TMtrans{\TMstroke}{\TMstroke}{\TMright}} (B)
            edge node {\TMtrans{\TMblank}{\TMblank}{\TMleft}} (C)
        (C) edge [loop above] node {\TMtrans{\TMstroke}{\TMblank}{\TMstay}} (C);
\end{tikzpicture}
\]
Alih-alih berhenti dalam keadaan~$q_2$, kita ingin melanjutkan operasi untuk
  menggandakan keluarannya. Ingatlah bahwa mesin pengganda menghapus simbol
  $\TMstroke$ pertama dalam masukan dan menuliskan dua simbol~$\TMstroke$ pada
  keluaran yang terpisah. Mari kita tambahkan instruksi untuk memastikan bahwa
  kepala pita membaca simbol~$\TMstroke$ pertama keluaran mesin penjumlahan.
\[
\begin{tikzpicture}[->,>=stealth',shorten >=1pt,auto,node distance=2.8cm,
                    semithick]
  \tikzstyle{every state}=[fill=none,draw=black,text=black]

  \node[initial,state] (A)              {$q_0$};
  \node[state]         (B) [right of=A] {$q_1$};
  \node[state]         (C) [right of=B] {$q_2$};
  \node[state]         (D) [below left of=C] {$q_3$};
  \node[state]         (E) [below left of=D] {$q_4$};

  \path (A) edge node {\TMtrans{\TMblank}{\TMstroke}{\TMstay}} (B)
            edge [loop above] node {\TMtrans{\TMstroke}{\TMstroke}{\TMright}} (A)
        (B) edge [loop above] node {\TMtrans{\TMstroke}{\TMstroke}{\TMright}} (B)
            edge node {\TMtrans{\TMblank}{\TMblank}{\TMleft}} (C)
        (C) edge node {\TMtrans{\TMstroke}{\TMblank}{\TMleft}} (D)
        (D) edge [loop left] node {\TMtrans{\TMstroke}{\TMstroke}{\TMleft}} (D)
            edge node[left, xshift=-2mm] {\TMtrans{\TMendtape}{\TMendtape}{\TMright}} (E);
\end{tikzpicture}
\]
Sekarang masukan itu mudah digandakan---kita hanya perlu menyambungkan mesin
pengganda pada keadaan~$q_4$. Hal ini mengharuskan kita mengganti nama keadaan
mesin pengganda agar dimulai dari~$q_4$, bukan~$q_0$; dengan demikian, kita
tidak memperoleh dua keadaan awal. Diagram akhirnya harus tampak seperti
dalam \olref{fig:combined}.
\begin{figure}
\[
\begin{tikzpicture}[->,>=stealth',shorten >=1pt,auto,node distance=2.8cm,
                    semithick]
  \tikzstyle{every state}=[fill=none,draw=black,text=black]
  \node[initial,state] (A)              {$q_0$};
  \node[state]         (B) [right of=A] {$q_1$};
  \node[state]         (C) [right of=B] {$q_2$};
  \node[state]         (D) [below left of=C] {$q_3$};
  \node[state]         (E) [below left of=D] {$q_4$};
  \node[state]         (2) [right of=E] {$q_5$};
  \node[state]         (3) [right of=2] {$q_6$};
  \node[state]         (4) [below of=3] {$q_7$};
  \node[state]         (5) [left of=4]  {$q_8$};
  \node[state]         (6) [left of=5]  {$q_9$};

  \path (A) edge node {\TMtrans{\TMblank}{\TMstroke}{\TMstay}} (B)
            edge [loop above] node {\TMtrans{\TMstroke}{\TMstroke}{\TMright}} (A)
        (B) edge [loop above] node {\TMtrans{\TMstroke}{\TMstroke}{\TMright}} (B)
            edge node {\TMtrans{\TMblank}{\TMblank}{\TMleft}} (C)
        (C) edge node {\TMtrans{\TMstroke}{\TMblank}{\TMleft}} (D)
        (D) edge [loop left] node {\TMtrans{\TMstroke}{\TMstroke}{\TMleft}} (D)
            edge node[left, xshift=-2mm] {\TMtrans{\TMendtape}{\TMendtape}{\TMright}} (E)
    (E) edge node {\TMtrans{\TMstroke}{\TMblank}{\TMright}} (2)
    (2) edge [loop above] node {\TMtrans{\TMstroke}{\TMstroke}{\TMright}} (2)
      edge node {\TMtrans{\TMblank}{\TMblank}{\TMright}} (3)
    (3) edge [loop above] node {\TMtrans{\TMstroke}{\TMstroke}{\TMright}} (3)
        edge node {\TMtrans{\TMblank}{\TMstroke}{\TMright}} (4)
    (4) edge [loop below] node {\TMtrans{\TMblank}{\TMstroke}{\TMleft}} (4)
        edge node {\TMtrans{\TMstroke}{\TMstroke}{\TMleft}} (5)
    (5) edge [loop below]  node {\TMtrans{\TMstroke}{\TMstroke}{\TMleft}} (5)
        edge              node {\TMtrans{\TMblank}{\TMblank}{\TMleft}} (6)
    (6) edge [loop below] node {\TMtrans{\TMstroke}{\TMstroke}{\TMleft}} (6)
        edge              node {\TMtrans{\TMblank}{\TMblank}{\TMright}} (E);
\end{tikzpicture}
\]
\caption{Menggabungkan mesin penjumlahan dan mesin pengganda}
\ollabel{fig:combined}
\end{figure}
\end{ex}

\begin{prop}
    Jika $M$ dan $M'$ bersifat disiplin dan masing-masing menghitung fungsi
    $f\colon \Nat^k \to \Nat$ dan $f'\colon \Nat \to \Nat$, maka
    $M \frown M'$ bersifat disiplin dan menghitung~$\comp{f}{f'}$.
\end{prop}

\begin{proof}
    Karena $M$ bersifat disiplin, ketika mesin itu berhenti dengan keluaran
    $f(n_1,\dots,n_k) = m$, kepalanya sedang memindai petak~$1$. Jika kita
    kemudian memasuki keadaan awal~$M'$, mesin akan berhenti dengan keluaran
    $f'(m)$ dan kembali memindai petak~$1$. Syarat-syarat lain dalam
    \olref[dis]{defn:disciplined} juga terpenuhi.
\end{proof}

\begin{prob}
    Berikan mesin Turing disiplin yang menghitung $f(x) = x+2$ dengan
    mengambil mesin~$M$ dari \cref{tur:mac:dis:prob:disc-succ} dan
    membentuk $M \frown M$.
\end{prob}
\end{document}
